\documentclass[11pt,a4paper]{article}
\usepackage[margin=24mm]{geometry}
\usepackage[T1]{fontenc}
\usepackage[utf8]{inputenc}
\usepackage{lmodern}
\usepackage{microtype}
\usepackage{graphicx}
\usepackage{booktabs}
\usepackage{tabularx}
\usepackage{array}
\usepackage{amsmath,amssymb}
\usepackage{float}
\usepackage{caption}
\usepackage{subcaption}
\usepackage{hyperref}
\usepackage{xcolor}
\usepackage{enumitem}
\usepackage{fancyhdr}
\usepackage{siunitx}
\sisetup{group-separator={,},group-minimum-digits=4}
\graphicspath{{lab1_report_assets/}}
\hypersetup{colorlinks=true,linkcolor=blue!55!black,urlcolor=blue!55!black}
\setlength{\parskip}{6pt}
\setlength{\parindent}{0pt}
\setlist[itemize]{topsep=3pt,itemsep=2pt}
\setlength{\headheight}{14pt}
\pagestyle{fancy}
\fancyhf{}
\lhead{ID2223 --- Lab 1}
\rhead{CIFAR-10: Tasks 1--5}
\cfoot{\thepage}
\newcommand{\figfull}[3]{\begin{figure}[htbp]\centering\includegraphics[width=#2\linewidth]{#1}\caption{#3}\end{figure}}
\newcommand{\pp}{\,\text{percentage points}}

\begin{document}
\begin{titlepage}
\centering
\vspace*{28mm}
{\Large KTH Royal Institute of Technology\par}
\vspace{8mm}
{\large ID2223 --- Scalable Machine Learning and Deep Learning\par}
\vspace{20mm}
{\Huge\bfseries Lab 1: CIFAR-10 Classification\par}
\vspace{6mm}
{\LARGE Tasks 1--5: Models, Regularization, Robustness and Adversarial Examples\par}
\vspace{20mm}
{\large Experimental report based on the supplied Jupyter notebooks\par}
\vfill
\begin{tabular}{rl}
Student(s): & \rule{85mm}{0.4pt}\\[5mm]
Date: & \today\\[5mm]
Implementation: & PyTorch / Google Colab\\
Data: & CIFAR-10\\
\end{tabular}
\vspace{25mm}
\end{titlepage}
\tableofcontents
\newpage

\section{Introduction and experimental scope}
This report studies image classification using CIFAR-10 with two convolutional neural networks: a CIFAR-adapted AlexNet and a CIFAR-adapted ResNet-18. The experiments comprise (1) dataset inspection and reproducible preparation, (2) baseline architecture and optimization comparisons, (3) controlled model-capacity and regularization studies, (4) corruption robustness with data augmentation, and (5) qualitative and adversarial prediction analysis. The report follows the concepts-first narrative of the accompanying practical ML foundations notebook: define the task, inspect the data, control the experiment, report measurements, and interpret their limitations.

The results and figures in this report are taken from the recorded outputs in the supplied combined Jupyter notebook and its source experiment notebook. They have \emph{not} been recomputed for this document. The digit-classification tutorial in Part A provided procedural background, while the reported CIFAR-10 results come from Tasks 1--5.

\section{Task 1 --- Study and prepare the data}
\subsection{What an example represents and class distribution}
CIFAR-10 contains \num{60000} color photographs, each of size $32\times32$ pixels and labelled as one of ten mutually exclusive classes: airplane, automobile, bird, cat, deer, dog, frog, horse, ship, and truck. A single example is an RGB tensor with shape $(3,32,32)$, containing \num{3072} channel intensity values. The original pixel representation is an unsigned 8-bit integer in $[0,255]$.

The provided dataset has \num{50000} training and \num{10000} official test images. Every class has \num{5000} examples in the provided training set and \num{1000} in the test set. Thus accuracy is not distorted by a major class-frequency imbalance, although per-class performance remains important.

\figfull{cifar_samples.png}{0.76}{Examples from each CIFAR-10 class, extracted from the notebook. Small resolution, background variation, and viewpoint changes make recognition challenging.}

\subsection{Training, validation and test partitioning}
The original \num{50000} training examples are divided by a stratified random split into \num{45000} training images and \num{5000} validation images, using \texttt{SPLIT\_SEED = 2223}. Stratification gives \num{4500} training and \num{500} validation examples per class. The official \num{10000} examples are held apart as a final test set. The task-specific corruption studies use the validation partition; the Task 5 classifier and FGSM analyses use the test partition.

The implementation retains sample IDs to trace predictions to their source images. Batches are read from reusable split objects, shuffled with a generator during training, and processed without shuffling during deterministic evaluation.

\subsection{Preprocessing for both architectures}
Both models consume normalized $(B,3,32,32)$ tensors. First, each raw image is converted from uint8 to float and divided by 255. The per-channel mean and standard deviation are estimated \emph{only} from the training partition:
\begin{equation}
\boldsymbol{\mu}=(0.4913,0.4822,0.4464),\quad
\boldsymbol{\sigma}=(0.2470,0.2434,0.2615).
\end{equation}
Then each pixel channel is standardized as $x'_{c,h,w}=(x_{c,h,w}-\mu_c)/\sigma_c$. The same values are used for training, validation, test and perturbed images. No model-specific resizing is required because both architectures are adapted to the original CIFAR resolution. Importantly, augmentation is applied in the $[0,1]$ pixel domain \emph{before} normalization, when enabled.

\subsection{Control of randomness and computational setting}
The notebook specifies \texttt{BASE\_SEED = 42} for initialization and experiment randomness, \texttt{SPLIT\_SEED = 2223} for partitions, and explicit generators for data order and stochastic transformations. It seeds Python, NumPy and PyTorch and configures deterministic behaviour where supported. Exact bitwise reproducibility across runtimes is not guaranteed: CUDA kernels, libraries, accelerator versions, and interrupted/resumed runs may differ. The recorded hardware is a Google Colab \textbf{Tesla T4 GPU}; an earlier ResNet epoch was executed on CPU before switching runtimes.

The selected \texttt{local} budget is eight epochs, batch size 64, full \num{45000}/\num{5000} training/validation partitions, and one principal seed. The experiments favor a small, interpretable set of controlled changes over a large parameter search.

\section{Task 2 --- Implement and compare AlexNet and ResNet-18}
\subsection{AlexNet adaptation, convolution operations and head}
The implemented AlexNet is a small-image adaptation, not a reproduction of the original ImageNet-scale architecture. Its five $3\times3$ convolutional layers have output channels $64,192,384,256,256$, with ReLU activations and three $2\times2$ stride-2 max-pooling operations. A convolution uses shared learned filters over local image neighborhoods; early layers can respond to edges, textures and colors, while deeper layers combine these responses into more complex visual patterns. Pooling reduces spatial dimensions and increases the effective receptive field.

\begin{table}[htbp]\centering\small
\caption{CIFAR-10 AlexNet feature-map and classifier dimensions.}
\begin{tabular}{ll}\toprule
Operation & Output per image\\\midrule
Input & $3\times32\times32$\\
Conv 1, ReLU, pool & $64\times16\times16$\\
Conv 2, ReLU, pool & $192\times8\times8$\\
Conv 3, ReLU & $384\times8\times8$\\
Conv 4, ReLU & $256\times8\times8$\\
Conv 5, ReLU, pool & $256\times4\times4$\\
Flatten & $4096$\\
Classifier & $4096\rightarrow1024\rightarrow512\rightarrow10$\\\bottomrule
\end{tabular}\end{table}

The classifier has two ReLU hidden layers, each followed by dropout with probability 0.5, and an output layer producing ten logits. The model has \textbf{\num{6976842} trainable parameters}. Cross-entropy computes the multi-class objective directly from logits; softmax is applied when probabilities are required for interpretation.

\subsection{ResNet-18 adaptation and residual connections}
The implemented ResNet-18 begins with a CIFAR-sized $3\times3$ convolution rather than the original larger stem. It uses four residual stages with two basic blocks each. Every basic block has two $3\times3$ convolutions with batch normalization and ReLU, plus an additive shortcut. When dimensions match, the shortcut is an identity map; when resolution or channel count changes, a stride-adjusted $1\times1$ convolution and batch normalization project the shortcut to the required shape.

\begin{equation}
\mathbf{y}=\operatorname{ReLU}\!\left(F(\mathbf{x})+S(\mathbf{x})\right),
\end{equation}
where $F$ is the learned residual transformation and $S$ is the identity or projection shortcut. Shortcuts help propagate information and gradients through deeper networks and make residual transformations easier to optimize.

The spatial/channel stages are $64\times32\times32$, $128\times16\times16$, $256\times8\times8$ and $512\times4\times4$. Global average pooling reduces the last tensor to 512 values and a linear layer maps these values to ten logits. The model has \textbf{\num{11173962} trainable parameters}. Both networks passed a dummy-batch test producing output shape $(2,10)$.

\subsection{Training, checkpointing and performance}
The baseline setup uses SGD with learning rate 0.01, momentum 0.9, weight decay $5\times10^{-4}$, batch size 64 and eight epochs. The notebook records per-epoch training/validation loss and accuracy. The reusable experiment routine retains the most recent checkpoint, selects its \emph{best} checkpoint by validation accuracy, and saves model, optimizer, and generator state where supported. Baseline AlexNet was also trained through a manual epoch loop with a saved weight-only checkpoint.

\begin{table}[htbp]\centering
\caption{Task 2 baseline results. Validation figures are taken from the recorded experiment; test figures are from the official test split.}
\begin{tabular}{lrrrr}\toprule
Model & Parameters & Best val. acc. & Test acc. & Test loss\\\midrule
AlexNet & \num{6976842} & 80.20\% & 78.71\% & 0.6483\\
ResNet-18 & \num{11173962} & 80.98\% & 81.00\% & 0.6650\\\bottomrule
\end{tabular}\end{table}

\textbf{Interpretation.} ResNet-18 gives higher test accuracy by 2.29 percentage points, at the cost of about 1.60 times as many trainable parameters. Validation and test losses should not be interpreted interchangeably, nor do small metric differences establish statistical significance with one seed. The recorded ResNet training duration included an approximately \num{4052.7}-second CPU epoch followed by GPU execution, so its total is \emph{not} a valid GPU-only timing comparison with AlexNet.

\section{Task 3 --- Investigate capacity and regularization}
All comparisons in this section use the same CIFAR-10 split, eight-epoch budget, SGD learning rate 0.01, momentum 0.9, batch size 64, and principal random seed. Only the named experimental variable is changed within each comparison. The shared baseline in Task 3 uses dropout 0.5 and weight decay 0.0005 and was separately logged with the reusable training routine; it should not be confused with the independently recorded Task 2 AlexNet baseline accuracy.

\subsection{Model capacity: classifier width}
The five-layer convolutional feature extractor is held fixed. Only the two hidden linear-layer widths change: small $(256,128)$, baseline $(1024,512)$, and large $(2048,1024)$.

\begin{table}[htbp]\centering\small
\caption{Capacity comparison: measured eight-epoch outcomes.}
\begin{tabular}{lrrrr}\toprule
Configuration & Parameters & Final train acc. & Best val. acc. & Time (s)\\\midrule
Small & \num{3334602} & 77.20\% & 76.72\% & 68.64\\
Baseline & \num{6976842} & 82.22\% & 78.56\% & 80.15\\
Large & \num{12750666} & 83.42\% & 79.06\% & 89.15\\\bottomrule
\end{tabular}\end{table}

\figfull{capacity_accuracy.png}{0.95}{Training and validation accuracy across eight epochs for three classifier widths.}
\figfull{capacity_loss.png}{0.95}{Training and validation cross-entropy loss for the capacity experiment.}

Increasing classifier width improves best validation accuracy from 76.72\% to 79.06\%, but the incremental gain shrinks: +1.84 percentage points from small to baseline, and +0.50 from baseline to large. Total training time increases from 68.64 to 89.15 seconds. The larger models show a wider end-of-run training/validation gap, but validation loss still generally decreases; this is insufficient to claim severe overfitting. The small model's poorer fit may reflect limited capacity and/or slower optimization under the fixed eight-epoch budget.

\subsection{L2 regularization through SGD weight decay}
In PyTorch SGD, weight decay corresponds to an L2-style penalty in the parameter update. The same baseline architecture and dropout $p=0.5$ are trained with weight decay $0$, $0.0005$ and $0.005$.
\begin{table}[htbp]\centering\small
\caption{L2 comparison, using the shared Task 3 baseline.}
\begin{tabular}{lrrrr}\toprule
Weight decay & Final train acc. & Best val. acc. & Best epoch & Time (s)\\\midrule
0 & 83.92\% & 79.94\% & 8 & 75.43\\
0.0005 & 82.22\% & 78.56\% & 8 & 80.15\\
0.005 & 67.64\% & 69.28\% & 8 & 79.12\\\bottomrule
\end{tabular}\end{table}
\figfull{l2_curves.png}{0.94}{Training and validation accuracy/loss under three L2-style weight-decay settings.}

No weight decay performed best after eight epochs (79.94\% validation accuracy). The strongest setting sharply lowers both training and validation accuracy, consistent with overly restrictive regularization or slower optimization at the chosen training duration. This comparison does not demonstrate that removing L2 is universally optimal; longer training and different random seeds could change the ranking.

\subsection{Dropout}
With architecture and weight decay fixed, dropout probabilities $p=0$, $0.3$ and $0.5$ are compared. Dropout temporarily zeros hidden activations during training; evaluation uses the full network.
\begin{table}[htbp]\centering\small
\caption{Dropout comparison. Training accuracy is from the final epoch; best validation accuracy may occur earlier.}
\begin{tabular}{lrrrr}\toprule
Dropout $p$ & Final train acc. & Best val. acc. & Best epoch & Time (s)\\\midrule
0.0 & 85.37\% & 78.64\% & 7 & 77.80\\
0.3 & 83.53\% & 78.46\% & 7 & 76.75\\
0.5 & 82.22\% & 78.56\% & 8 & 80.15\\\bottomrule
\end{tabular}\end{table}
\figfull{dropout_curves.png}{0.94}{Accuracy and loss trajectories for dropout probabilities 0, 0.3 and 0.5.}

Without dropout the network achieves its highest training accuracy and best validation accuracy at epoch 7, but validation accuracy decreases at epoch 8 (78.64\% to 77.64\%). This is a sign of emerging overfitting. Dropout reduces training accuracy, and all three best validation accuracies are within 0.18 percentage points, so the evidence for a particular optimal dropout level is weak over just eight epochs.

\subsection{Computational cost and control}
The capacity comparison is particularly informative for cost: 3.33, 6.98 and 12.75 million parameters correspond to 68.64, 80.15 and 89.15 seconds of measured eight-epoch training, respectively. Reported peak CUDA memory for those runs was approximately 396, 466 and 584 MB; however, other model objects can remain allocated in the Colab process. These are observed allocation peaks, \emph{not isolated per-model footprints}. For L2 and dropout, parameter counts remain unchanged at \num{6976842}; measured times are similar, and interpretation should prioritize accuracy differences rather than small runtime fluctuations.

\section{Task 4 --- Evaluate robustness}
\subsection{Training augmentation and controlled corruption}
Three AlexNet configurations are compared with the same architecture, learning rate, momentum, weight decay, training split and eight-epoch budget. The \emph{original} condition has no augmentation. \emph{Standard} augmentation applies reflective padding, random cropping and horizontal flipping. \emph{Strong} augmentation additionally includes randomly applied rotations up to 15 degrees, shears up to 10 degrees and additive Gaussian noise with scale sampled up to 0.06. These transformations are chosen to approximate modest viewpoint changes and sensor noise without routinely destroying class identity.

\figfull{augmentation_examples.png}{0.77}{Unmodified, standard-augmented, and strong-augmented examples from the notebook.}

For evaluation, the same \num{5000} validation images and labels are used in all three conditions: clean pixels, a fixed 15-degree rotation, and independent zero-mean Gaussian noise with standard deviation 0.06 on the $[0,1]$ scale. The noise generator uses a fixed seed of 42. Every model receives exactly the same transformed inputs; only its learned parameters differ. No gradient or optimizer update is performed on this evaluation set.

\subsection{Robustness measurements}
\begin{table}[htbp]\centering\small
\caption{Task 4 validation accuracy on the same 5,000 images under fixed corruption.}
\begin{tabular}{lrrr}\toprule
Training condition & Clean & Rotation $15^\circ$ & Gaussian noise\\\midrule
No augmentation & \textbf{78.56\%} & 70.32\% & 66.38\%\\
Standard crop/flip & 77.56\% & 68.44\% & 65.06\%\\
Strong augmentation & 76.40\% & \textbf{71.86\%} & \textbf{73.74\%}\\\bottomrule
\end{tabular}\end{table}
\figfull{robustness.png}{0.87}{Comparison of clean and corrupted validation accuracy for all three training conditions.}

Stronger augmentation costs 2.16 percentage points of clean accuracy relative to no augmentation, but improves rotated accuracy by 1.54 points and noisy accuracy by 7.36 points. In contrast, standard crop/flip alone does not improve robustness against these \emph{different} evaluation distortions. The accuracy loss from clean to noisy images is 12.18 points without augmentation versus only 2.66 points with strong augmentation.

\subsection{Newly introduced errors and failure modes}
A \emph{new error} is an example correctly classified in the clean condition and incorrectly classified after corruption. This quantity differs from the net accuracy loss, since some initially incorrect predictions can become correct after transformation.
\begin{table}[htbp]\centering
\caption{Corruption-induced errors among the same 5,000 validation IDs.}
\begin{tabular}{lrr}\toprule
Corruption & Original: new errors & Strong aug.: new errors\\\midrule
Rotation $15^\circ$ & 647 (12.94\%) & 443 (8.86\%)\\
Gaussian noise & 732 (14.64\%) & 252 (5.04\%)\\\bottomrule
\end{tabular}\end{table}
\figfull{corruption_examples.png}{0.90}{Examples correctly identified by the original AlexNet when clean, but misclassified after Gaussian noise.}

The displayed noise-induced errors include confusions between vehicle categories (truck vs. automobile) as well as unrelated class predictions (some vehicles or animals classified as frog). These examples show that corruption can disrupt both shape and texture cues. Stronger transformations are not always label preserving: at $32\times32$ resolution, larger rotations, crops or shears may remove a defining object part or make a borderline image ambiguous. The present results therefore apply to the specified moderate transformations, not arbitrary distortion strengths.

\subsection{Training and evaluation overhead}
Training with augmentation requires extra crop/flip and, for the strong setting, affine resampling and noise generation for every batch. The notebook records per-epoch times for these runs: the standard condition is roughly 10 seconds per epoch, and the strong condition roughly 11--12 seconds per epoch on the logged GPU. Evaluation of each model-condition pair on 5,000 images took on the order of a few tenths of a second in the reported runs, but those timings include execution variability and are not a dedicated inference benchmark. The larger GPU-memory peaks reported in later sequential experiments are confounded by other live model allocations and should not be attributed entirely to augmentation.

\section{Task 5 --- Analyse predictions and adversarial examples}
\subsection{Quantitative classification results}
The Task 2 AlexNet and ResNet-18 checkpoints are evaluated on the official \num{10000}-image test set. A classification report gives per-class precision, recall, and F1-score; confusion matrices identify systematic confusions.
\begin{table}[htbp]\centering
\caption{Whole-test performance, as reported by the notebook.}
\begin{tabular}{lrrrr}\toprule
Model & Accuracy & Macro precision & Macro recall & Macro F1\\\midrule
AlexNet & 78.7\% & 0.791 & 0.787 & 0.786\\
ResNet-18 & 81.0\% & 0.826 & 0.810 & 0.812\\\bottomrule
\end{tabular}\end{table}
\figfull{confusion_matrices.png}{0.99}{Confusion matrices for the two trained models on the official CIFAR-10 test set.}

Dog is particularly difficult: recall is 62.0\% for AlexNet and 59.5\% for ResNet-18. ResNet-18 misclassifies 267 actual dog images as cat, compared with 182 for AlexNet, although it improves cat recall from 62.3\% to 81.0\%. Other confusion patterns involve related animals and vehicle classes. This demonstrates why overall accuracy alone does not convey class-specific tradeoffs.

\subsection{Successful and unsuccessful predictions}
The experiment inspects four groups: both models correct, both wrong, only AlexNet correct, and only ResNet-18 correct. This avoids the selection bias of showing only cases where one model fails. Individual examples reveal challenging viewpoints, visually similar animals, low resolution and distracting backgrounds. Some incorrect predictions have high softmax confidence, emphasizing that confidence is not equivalent to correctness or probability calibration.
\figfull{prediction_categories.png}{0.95}{Examples organized by whether neither, one, or both models make a correct prediction. Text above each image records true class, predicted class and reported confidence.}

\subsection{FGSM attack, constraint and evaluation protocol}
The fast gradient sign method (FGSM) is a one-step, untargeted, white-box perturbation. For a frozen classifier with cross-entropy loss $\mathcal{L}$, its adversarial input is
\begin{equation}
\mathbf{x}_{\mathrm{adv}}=\operatorname{clip}_{[0,1]}\Bigl(\mathbf{x}+\varepsilon\operatorname{sign}\bigl(\nabla_{\mathbf{x}}\mathcal{L}(f_\theta(\mathbf{x}),y)\bigr)\Bigr).
\end{equation}
All gradients are taken with respect to \emph{input pixels}, not model weights. Input gradients pass through the training-set normalization. Clipping keeps pixel values valid, and the $L_\infty$ constraint is $\|\mathbf{x}_{\rm adv}-\mathbf{x}\|_\infty\leq\varepsilon$ for pixels on $[0,1]$. A verification batch recorded a maximum absolute perturbation of approximately 0.03 at $\varepsilon=0.03$, consistent with the constraint up to floating-point precision.

Both models were attacked separately on the \emph{same first 500 test images}, for $\varepsilon\in\{0,0.01,0.03,0.05\}$. Because each attack uses the gradients of its target model, the generated perturbations are not the same pixel arrays across models. This is an appropriate model-specific white-box robustness comparison rather than a cross-model transfer test.
\begin{table}[htbp]\centering
\caption{Accuracy on 500 test images under model-specific FGSM.}
\begin{tabular}{lrrrr}\toprule
Model & $\varepsilon=0$ & $0.01$ & $0.03$ & $0.05$\\\midrule
AlexNet & 76.4\% & 37.0\% & 6.8\% & 2.0\%\\
ResNet-18 & 81.0\% & 14.2\% & 1.8\% & 0.8\%\\\bottomrule
\end{tabular}\end{table}
\figfull{fgsm_examples.png}{0.90}{The same originally correct AlexNet images under increasing FGSM perturbation strength; predictions and confidence may change markedly despite modest visual alterations.}

Both models are vulnerable to deliberately optimized perturbations. In this specific one-step attack, ResNet-18 has higher clean accuracy but lower accuracy after perturbation. This does not show that AlexNet is universally more adversarially robust: attack type, optimization strength, model-specific gradients, seeds, image selection and defense settings matter. The $\varepsilon=0$ accuracy for AlexNet (76.4\%) is measured on 500 images, not the full test set, so it need not equal the 78.7\% whole-test accuracy.

\subsection{Random corruption versus adversarial perturbation}
Task 4 noise was sampled from a fixed stochastic distribution without consulting the model. FGSM instead uses access to the trained model's loss gradient to maximize error under a per-pixel constraint. Therefore a small FGSM perturbation can cause far larger accuracy loss than random noise of a superficially similar magnitude. Comparing them requires stating both the corruption distribution and the adversary's knowledge of the model.

\section{Discussion, limitations and conclusions}
Across Tasks 1--5, the reproducible CIFAR-10 split and shared preprocessing allow meaningful controlled comparisons. ResNet-18 attains higher overall test accuracy than the adapted AlexNet (81.0\% versus 78.7\%), although both have class-specific weaknesses. Increasing AlexNet classifier width gives modest improvements with greater training cost; excessively strong L2-style weight decay impedes fitting over eight epochs; dropout changes training dynamics but produces similar best validation accuracy in the short-run comparisons.

Robustness depends on the perturbation: strong training augmentation sacrifices some clean accuracy but reduces errors from fixed rotation and additive noise. Neither architecture is robust to the tested model-specific white-box FGSM attacks, demonstrating that ordinary corruption robustness and adversarial robustness are different properties.

\textbf{Limitations.} Most configurations were run with one seed and eight epochs. The runs include different checkpoint histories and at least one CPU-to-GPU transition, making aggregate training durations between Task 2 architectures noncomparable. Recorded CUDA memory can include multiple live models, and short evaluation times are noisy. Checkpoint selection uses validation accuracy, and selected configurations were inspected on the official test set only at the analysis stage. The Task 5 FGSM sample is only 500 images and uses a single one-step attack. Stronger image transformations could invalidate labels, while subtle distortions may affect individual classes differently.

\textbf{Useful follow-up.} A high-value next experiment would repeat the most promising configurations across multiple seeds with a consistent GPU-only training budget and report means and variability. A separate inference benchmark and a stronger iterative adversarial attack would help determine whether the observed clean-accuracy and robustness rankings are stable.

\section*{References and provenance}
\addcontentsline{toc}{section}{References and provenance}
\begin{enumerate}[leftmargin=*]
\item Course notebooks supplied for this experiment: \texttt{ID2223\_Lab1.ipynb}, \texttt{Lab1\_2\_5.ipynb}, and the combined \texttt{ID2223\_Lab1\_Combined\_Tasks1-5.ipynb}. Tables and included figures reproduce the supplied saved results.
\item A. Krizhevsky, \emph{Learning Multiple Layers of Features from Tiny Images}, 2009 (CIFAR-10 dataset).
\item A. Krizhevsky, I. Sutskever and G. Hinton, \emph{ImageNet Classification with Deep Convolutional Neural Networks}, 2012 (original AlexNet).
\item K. He, X. Zhang, S. Ren and J. Sun, \emph{Deep Residual Learning for Image Recognition}, 2016 (ResNet).
\item I. Goodfellow, J. Shlens and C. Szegedy, \emph{Explaining and Harnessing Adversarial Examples}, 2015 (FGSM).
\end{enumerate}
\end{document}
